\startchapter[title={Soft Sensor}]
	Definido a partir do inglês {\emph software sensor}. Tem o papel de substituir um sensor físico ou atuar como um complemento para este.

	\startplacefigure[title={Esquemático de um fermentador com seu aparato de medição e controle}, reference={fig:bioreactor_control}]
		\externalfigure[image/bioreactor_control.png][height=200pt]
	\stopplacefigure

	Um bioreator com variáveis medidas como o da \in{Figura}[fig:bioreactor_control] pode ter sua montagem física simplificada, barateada, e estendida.

	Pode ser utilizado para predizer falhas sistêmicas, {\emph fault detection}.
\stopchapter

\startchapter[title={Séries Temporais}]
	Séries temporais ou {\emph time series} capturam o comportamento dinâmico do sistema. Contrariamente métodos que levam em consideração somente o estado inicial, ou atual. A \in{Equação}[eq:time_series_inequality] exemplifica a ideia de um sistema variante no tempo.

	\startplaceformula[reference={eq:time_series_inequality}]
		\startformula
			f(x_{1}, x_{2}, \cdots, x_{n}, t_{0}) \neq \f(x_{1}, x_{2}, \cdots, x_{n}, t_{1})
		\stopformula
	\stopplaceformula

	Um sistema fermentativo é mais do que isso, ele depende dos estados anteriores também. Os físicos chamam esse tipo de sistema não markoviano.

	\startplacefigure[title={Exemplo de comportamento de uma variável em uma série temporal}, reference={fig:time_series_graph}]
		\externalfigure[image/time_series_graph.png][height=200pt]
	\stopplacefigure
\stopchapter

\startchapter[title={Fermentação}]
\stopchapter
